% =====================================================================
% Strukturplan für die Bachelorarbeit
% Eigenständiges Dokument (nicht Teil von thesis.tex) - dient als
% Planungs-/Referenzdokument und als erste Uebungsdatei fuer VS Code.
%
% Revision 2: eingearbeitet Rückmeldung des Betreuers vom 2026-08-13
% (siehe "Änderungen in dieser Revision" nach der Executive Summary).
% =====================================================================
\documentclass[a4paper,11pt]{article}

\usepackage[english]{babel}
\usepackage[T1]{fontenc}
\usepackage[utf8]{inputenc}

\usepackage{geometry}
\geometry{a4paper, left=25mm, right=25mm, top=25mm, bottom=25mm}

\usepackage{enumitem}
\setlist{itemsep=1pt, topsep=2pt}

\PassOptionsToPackage{hyphens}{url}
\usepackage[colorlinks=true, linkcolor=blue, urlcolor=blue, citecolor=blue]{hyperref}

\usepackage{xcolor}
\definecolor{thiblau}{rgb}{0,0.112,0.47}
\definecolor{notebg}{RGB}{240,240,250}
\definecolor{profbg}{RGB}{255,244,224}

\usepackage{titlesec}
\titleformat{\section}{\Large\bfseries\color{thiblau}}{\thesection}{1em}{}
\titleformat{\subsection}{\large\bfseries\color{thiblau}}{\thesubsection}{1em}{}

\usepackage{tcolorbox}
\tcbuselibrary{skins,breakable}

\title{\color{thiblau}Strukturplan für die Bachelorarbeit}
\author{Thomas Schneider}
\date{Revision 2 -- Stand: \today}

\begin{document}

\maketitle

% ---------------------------------------------------------------
\section*{General Information}
\begin{tabular}{@{}ll}
	\textbf{Name:}      & Thomas Schneider \\
	\textbf{Standort:}  & Technische Hochschule Ingolstadt (THI) \\
	\textbf{Sprache:}   & Englisch \\
	\textbf{Titel:}     & Implementation and Visualization of a Post-Quantum Cryptosystem \\
	                    & Using the Example of the McEliece Cryptosystem \\
\end{tabular}

% ---------------------------------------------------------------
\section*{Research Question}
How can the McEliece cryptosystem be implemented and visualized within CrypTool
Online to improve the understanding of code-based post-quantum cryptography for
students?

% ---------------------------------------------------------------
\section*{Executive Summary}
The emergence of quantum computing poses a significant threat to currently deployed
public-key cryptographic systems such as RSA and Elliptic Curve Cryptography.
Algorithms like Shor's algorithm are expected to compromise the security assumptions
of these systems once sufficiently powerful quantum computers become available.
Consequently, post-quantum cryptography has become an important research and
development area within modern information security.

Among the various approaches to post-quantum cryptography, the McEliece cryptosystem
represents one of the oldest and most extensively studied code-based cryptographic
systems. Despite decades of cryptanalytic research, no practical attacks against
properly parameterized variants of the system are known. However, due to its strong
mathematical foundations in coding theory, the McEliece cryptosystem is often
perceived as difficult to understand, especially for students encountering
post-quantum cryptography for the first time.

This thesis addresses the challenge of making the McEliece cryptosystem more
accessible for educational purposes. Rather than focusing on advanced mathematical
proofs or cryptanalytic research, the thesis investigates how the essential concepts
and processes of the cryptosystem can be presented in a comprehensible and
interactive manner.

The primary objective is the design and implementation of an educational
visualization tool based on CrypTool Online (CTO). The work includes the
identification of relevant theoretical concepts, the development of a didactic
visualization approach, the design of a software architecture suitable for
educational use, and the implementation of a CrypTool Online plugin that visualizes
the key generation, encryption, and decryption processes of the McEliece
cryptosystem.

The resulting artifact aims to support students in understanding the relationship
between coding theory and post-quantum cryptography while demonstrating how
interactive visualizations can improve the learning process in complex
security-related domains.

The thesis contributes both a practical software implementation and a didactic
framework for teaching code-based post-quantum cryptography.

% ---------------------------------------------------------------
\clearpage
\section*{Änderungen in dieser Revision}

\begin{tcolorbox}[breakable, colback=profbg, colframe=thiblau, title=Rückmeldung des Betreuers (2026-08-13)]
\begin{enumerate}
	\item \textbf{Zu 3.1--3.3 (Quantum Computing, Shor's, Grover's Algorithm):}
	      ,,Sie können überlegen diese Punkte zu einem Grundlagenkapitel
	      zusammenzufassen.``
	      $\rightarrow$ zu \emph{einer} Subsection zusammengefasst (jetzt 3.1).

	\item \textbf{Zu Kapitel 5 \& 6 (CrypTool Online Architecture / Educational
	      and Didactic Concept):} ,,6 würde ich vor 5 ziehen, dann haben Sie
	      alle Grundlagen hintereinander.``
	      $\rightarrow$ Reihenfolge getauscht: CrypTool Online Architecture ist
	      jetzt Kapitel 5, Educational and Didactic Concept Kapitel 6.

	\item \textbf{Zu 1.2 (Problem Statement):} ,,Hier kürzen -- die Einleitung
	      sollte maximal 2 Seiten haben, lieber schieben Sie einige Teile
	      (Problem Statement) weiter hinten, nach den Grundlagen in ein
	      'Requirementskapitel'. Da können Sie dann auch die Objectives
	      ausarbeiten und den Scope nochmal herausarbeiten.``
	      $\rightarrow$ Problem Statement, Objectives und Scope aus der
	      Einleitung entfernt; neues Kapitel~7 \emph{Requirements} nach den
	      Grundlagenkapiteln eingefügt, das diese drei Punkte aufnimmt und
	      ausarbeitet.

	\item \textbf{Zu Kapitel 2 (Fundamentals of Cryptography):} ,,Hier dürfen
	      Sie knapp bleiben: Fokus, was für die Arbeit wichtig ist. Zielgruppe
	      Ihrer Arbeit sind Absolventen Ihres Studiengangs, d.\,h. einiges
	      wissen die auch.``
	      $\rightarrow$ Gliederung unverändert gelassen (4 Subsections), aber
	      als Hinweis direkt im Kapitel vermerkt: knapp halten, nur worauf
	      spätere Kapitel wirklich aufbauen.
\end{enumerate}
\end{tcolorbox}

\clearpage
\tableofcontents
\clearpage

% =================================================================
\section{Introduction}
\textit{Ziel: max. 2 Seiten.}

\subsection{Motivation}
\begin{itemize}
	\item Bedrohung klassischer Kryptographie durch Quantencomputer
	\item Relevanz von PQC
	\item Schwierigkeit des Verständnisses codebasierter Verfahren
	\item Bedarf an Lehr- und Lernwerkzeug
\end{itemize}

\subsection{Research Question}
Vorstellung der Forschungsfrage

\subsection{Structure of the Thesis}
Kurze Beschreibung aller Kapitel

% =================================================================
\section{Fundamentals of Cryptography}
\textit{Knapp halten -- Zielgruppe (Absolvent*innen des eigenen Studiengangs)
kennt die Grundlagen bereits. Fokus nur auf das, was für Kapitel 4
(McEliece) und die Abgrenzung zu klassischen Verfahren wirklich gebraucht
wird.}

\subsection{Symmetric and Asymmetric Cryptography}
\begin{itemize}
	\item Grundprinzipien
	\item Key Exchange
	\item Public-Key-Verfahren
\end{itemize}

\subsection{Security Concepts}
\begin{itemize}
	\item Confidentiality
	\item Integrity
	\item Authenticity
\end{itemize}

\subsection{Computational Hardness}
\begin{itemize}
	\item One-Way Functions
	\item Trapdoor Functions
\end{itemize}

\subsection{Public-Key Cryptography}
\begin{itemize}
	\item RSA
	\item ECC
	\item Typische Sicherheitsannahmen
\end{itemize}

% =================================================================
\section{Quantum Computing and Post-Quantum Cryptography}

\subsection{Quantum Computing Fundamentals and Impact on Classical Cryptography}
\textit{Zusammengefasst aus den ehemals separaten Punkten "Fundamentals of
Quantum Computing", "Shor's Algorithm" und "Grover's Algorithm" (Vorschlag
des Betreuers).}
\begin{itemize}
	\item Qubits, Superposition, Quantum Advantage
	\item Shor's Algorithm -- Auswirkungen auf RSA und ECC
	\item Grover's Algorithm -- Auswirkungen auf symmetrische Verfahren
\end{itemize}

\subsection{Need for Post-Quantum Cryptography}
\begin{itemize}
	\item Store Now, Decrypt Later
	\item Migration Challenges
\end{itemize}

\subsection{Overview of PQC Families}
\begin{itemize}
	\item Lattice-based
	\item Code-based
	\item Hash-based
	\item Multivariate
	\item Isogeny-based
\end{itemize}

\subsection{Position of McEliece within PQC}
\begin{itemize}
	\item Historische Entwicklung
	\item NIST-Prozess
	\item Classic McEliece
\end{itemize}

% =================================================================
\section{The McEliece Cryptosystem}

\subsection{Error-Correcting Codes}
\begin{itemize}
	\item Motivation
	\item Kommunikationsfehler
	\item Fehlerkorrektur
\end{itemize}

\subsection{Linear Codes}
\begin{itemize}
	\item Codewords
	\item Generator Matrix
	\item Parity Check Matrix
\end{itemize}

\subsection{Goppa Codes}
\begin{itemize}
	\item Grundidee
	\item Eigenschaften
	\item Fehlerkorrekturfähigkeit
\end{itemize}

\subsection{Structure of the McEliece Cryptosystem}
\begin{itemize}
	\item Private Key
	\item Public Key
	\item Message
	\item Error Vector
\end{itemize}

\subsection{Key Generation}
\begin{itemize}
	\item Generator Matrix
	\item Scrambling Matrix
	\item Permutation Matrix
	\item Public Key Construction
\end{itemize}

\subsection{Encryption}
\begin{itemize}
	\item Encoding
	\item Error Addition
\end{itemize}

\subsection{Decryption}
\begin{itemize}
	\item Decoding
	\item Error Correction
\end{itemize}

\subsection{Security Assumption}
\begin{itemize}
	\item General Decoding Problem
	\item Schwierigkeit zufälliger linearer Codes
\end{itemize}

\subsection{Educational Abstraction}
\ %

% =================================================================
\section{CrypTool Online Architecture}
\textit{Vorgezogen von ehemals Kapitel 6 (Rückmeldung des Betreuers), damit
alle Grundlagenkapitel (2--6) hintereinander stehen.}

\subsection{Introduction to CrypTool}
\begin{itemize}
	\item CrypTool Projekt
	\item CrypTool Online
\end{itemize}

\subsection{Architecture of CrypTool Online}
\begin{itemize}
	\item Workflow-Konzept
	\item Plugins
	\item Datenflüsse
\end{itemize}

\subsection{Existing Educational Components}
\begin{itemize}
	\item Analyse bestehender CTO-Plugins
\end{itemize}

% =================================================================
\section{Educational and Didactic Concept}
\textit{Ehemals Kapitel 5, jetzt nach CrypTool Online Architecture (Rückmeldung
des Betreuers).}

\subsection{Learning Objectives}
\begin{itemize}
	\item PQC verstehen
	\item McEliece einordnen können
	\item Schlüsselabläufe nachvollziehen können
\end{itemize}

\subsection{Target Audience Analysis}
\begin{itemize}
	\item Bachelor-Studierende
	\item Grundkenntnisse Kryptographie
\end{itemize}

\subsection{Educational Challenges}
\begin{itemize}
	\item Abstrakte Mathematik
	\item Hohe Komplexität
\end{itemize}

\subsection{Visualization Principles}
\begin{itemize}
	\item Cognitive Load Theory
	\item Progressive Disclosure
	\item Step-by-Step Visualization
\end{itemize}

\subsection{Concept of the Learning Tool}
\begin{itemize}
	\item Lernszenarien
	\item Benutzerfluss
	\item Interaktionen
\end{itemize}

% =================================================================
\section{Requirements}
\textit{Neues Kapitel (Rückmeldung des Betreuers). Nimmt Problem Statement,
Objectives und Scope aus der (jetzt kurzen) Einleitung auf und arbeitet sie
aus; ergänzt um die technischen Anforderungen, die vorher unter "Design"
bzw. "CrypTool Online Architecture" standen (eigener Vorschlag, siehe
Anmerkungen am Ende des Dokuments).}

\subsection{Problem Statement}
\begin{itemize}
	\item McEliece ist mathematisch anspruchsvoll
	\item Vorhandene Literatur richtet sich häufig an Experten
	\item Fehlende interaktive Lernwerkzeuge
\end{itemize}

\subsection{Objectives}
\begin{itemize}
	\item Entwicklung eines CTO-Plugins
	\item Visualisierung zentraler Abläufe
	\item Unterstützung des Lernprozesses
\end{itemize}

\subsection{Scope}
Abgrenzung der Arbeit

\subsection{Functional Requirements}
\ %

\subsection{Non-Functional Requirements}
\ %

\subsection{Requirements for the McEliece Plugin}
\begin{itemize}
	\item Fachliche Anforderungen
	\item Didaktische Anforderungen
	\item Technische Anforderungen
\end{itemize}

% =================================================================
\section{Design of the McEliece Learning Plugin}

\subsection{System Design}
\begin{itemize}
	\item Architekturdiagramm
\end{itemize}

\subsection{User Interface Design}
\begin{itemize}
	\item Layout
	\item Navigation
	\item Interaktion
\end{itemize}

\subsection{Visualization Design}
\begin{itemize}
	\item Prozessdiagramme
	\item Matrixdarstellung
	\item Animationen
\end{itemize}

\subsection{Data Flow Design}
\begin{itemize}
	\item Eingaben
	\item Ausgaben
	\item Zwischenzustände
\end{itemize}

% =================================================================
\section{Implementation}

\subsection{Development Environment}
\begin{itemize}
	\item CTO Framework
	\item Programmiersprache
	\item Tools
\end{itemize}

\subsection{Implementation of Key Generation}
\begin{itemize}
	\item Technische Umsetzung
\end{itemize}

\subsection{Implementation of Encryption}
\begin{itemize}
	\item Technische Umsetzung
\end{itemize}

\subsection{Implementation of Decryption}
\begin{itemize}
	\item Technische Umsetzung
\end{itemize}

\subsection{Visualization Components}
\begin{itemize}
	\item Animationen
	\item Grafiken
	\item Interaktive Elemente
\end{itemize}

\subsection{Challenges and Solutions}
\begin{itemize}
	\item Technische Probleme
	\item Didaktische Probleme
\end{itemize}

% =================================================================
\section{Evaluation}
% siehe eigene Anmerkung am Ende des Dokuments zu moeglichen Unterkapiteln

% =================================================================
\section{Conclusion and Future Work}
% siehe eigene Anmerkung am Ende des Dokuments zu moeglichen Unterkapiteln

% =================================================================
\section{Sources}
\begin{itemize}
	\item Felix Peter Paul: \emph{Codebasierte Post-Quanten-Kryptografie: Goppa
	      Codes und das McEliece Kryptosystem}. Springer Vieweg, BestMasters,
	      2025. \url{https://doi.org/10.1007/978-3-658-46743-2} \\
	      (liegt als PDF im Ordner \texttt{Quellen/})

	\item Seng W. Loke: \emph{From Distributed Quantum Computing to Quantum
	      Internet Computing: An Introduction}. Wiley-IEEE Press, 2024. \\
	      (liegt als PDF im Ordner \texttt{Quellen/})

	\item Bundesamt für Sicherheit in der Informationstechnik (BSI):
	      Quantentechnologien und Post-Quanten-Kryptografie\\
	      \url{https://www.bsi.bund.de/EN/Themen/Unternehmen-und-Organisationen/Informationen-und-Empfehlungen/Quantentechnologien-und-Post-Quanten-Kryptografie/quantentechnologien-und-post-quanten-kryptografie_node.html}

	\item Classic McEliece (Web-Archiv-Version vom 25.12.2025)\\
	      \url{https://web.archive.org/web/20251225181518/https://classic.mceliece.org/index.html}

	\item BSI: Kryptografie quantensicher gestalten\\
	      \url{https://www.bsi.bund.de/SharedDocs/Downloads/DE/BSI/Publikationen/Broschueren/Kryptografie-quantensicher-gestalten.pdf?__blob=publicationFile&v=6}

	\item R.\ McEliece. \emph{A public-key cryptosystem based on algebraic coding
	      theory}. The Deep Space Network Progress Report, DSN PR 42--44, 1978.

	\item \textit{Neu vorgeschlagen (siehe Anmerkungen unten):} Bernstein,
	      Buchmann, Dahmen (Hrsg.), \emph{Post-Quantum Cryptography}, Springer,
	      2009; Shor (1994); Grover (1996); NIST PQC Project.
\end{itemize}
Alle Einträge sind bereits als Vorlage in \texttt{literature.bib} angelegt und
können per \verb|\cite{...}| referenziert werden.

% =================================================================
\clearpage
\section*{Eigene Anmerkungen und Verbesserungsvorschläge}
\textit{Die folgenden Punkte sind \textbf{keine} Vorgaben des Professors, sondern
eigene Vorschläge, die zur Diskussion stehen.}

\begin{tcolorbox}[breakable, colback=notebg, colframe=thiblau, title=Vorschläge]
\begin{enumerate}
	\item \textbf{Requirements für das Plugin (6.4 alt) in Kapitel 7 verschoben.}
	      Damit stehen alle Anforderungen (Problem Statement, Objectives, Scope,
	      Functional/Non-Functional Requirements, Requirements for the Plugin)
	      gebündelt an einer Stelle statt über zwei Kapitel verteilt.

	\item \textbf{Kapitel 10 (Evaluation) ausformulieren.} Aktuell hat dieses
	      Kapitel keine Unterkapitel. Vorschlag für eine Struktur:
	      \begin{itemize}
	      	\item 10.1 Evaluation Method (z.\,B. Experten-Review, kleine
	      	      Nutzerstudie mit Kommilitonen, Vergleich mit Lernzielen aus 6.1)
	      	\item 10.2 Study Design / Participants
	      	\item 10.3 Results
	      	\item 10.4 Discussion of Results
	      	\item 10.5 Limitations of the Evaluation
	      \end{itemize}

	\item \textbf{Kapitel 11 (Conclusion and Future Work) ebenfalls untergliedern:}
	      \begin{itemize}
	      	\item 11.1 Summary of Contributions
	      	\item 11.2 Answering the Research Question
	      	\item 11.3 Limitations
	      	\item 11.4 Future Work
	      \end{itemize}

	\item \textbf{Related Work ergänzen.} In Kapitel 1 oder 5.3 fehlt noch ein
	      expliziter Vergleich mit bereits existierenden Lehrmitteln zu
	      code-basierter/Post-Quanten-Kryptographie (z.\,B. andere
	      CrypTool-Plugins, interaktive Tools anderer Universitäten). Das stärkt
	      die Argumentation, warum ein neues Plugin nötig ist.

	\item \textbf{Abgrenzung von "Implementation" und "Visualisierung" schärfen.}
	      Der Titel nennt beides gleichwertig; in der Gliederung dominiert aktuell
	      die Implementierung (Kapitel 4, 8, 9). Es könnte helfen, in 7.3 (Scope)
	      explizit festzulegen, welche McEliece-Parametergrößen real unterstützt
	      werden (Lehr-/Demo-Parameter vs. reale/klassische Sicherheitsparameter),
	      da vollständige Classic-McEliece-Parameter für eine Visualisierung im
	      Browser vermutlich zu groß sind.

	\item \textbf{Zusätzliche Quellenvorschläge} (noch nicht offiziell Teil der
	      Quellenliste, aber etablierte Standardreferenzen):
	      \begin{itemize}
	      	\item D.\,J. Bernstein, J. Buchmann, E. Dahmen (Hrsg.),
	      	      \emph{Post-Quantum Cryptography}, Springer, 2009 -- gutes
	      	      Überblickswerk für Kapitel 3.
	      	\item P.\,W. Shor, \emph{Algorithms for quantum computation: discrete
	      	      logarithms and factoring}, FOCS 1994 -- Primärquelle für 3.1.
	      	\item L.\,K. Grover, \emph{A fast quantum mechanical algorithm for
	      	      database search}, STOC 1996 -- Primärquelle für 3.1.
	      	\item NIST Post-Quantum Cryptography Project
	      	      (\url{https://csrc.nist.gov/projects/post-quantum-cryptography})
	      	      -- für 3.4 (NIST-Prozess, Classic McEliece Status).
	      \end{itemize}
	      Bitte gegenprüfen/ergänzen -- ich habe diese nicht aus Ihrem
	      Quellen-Ordner, sondern aus meinem Wissen vorgeschlagen.

	\item \textbf{Für Kapitel 5/9 (CrypTool Online):} Sobald es inhaltlich
	      sinnvoll ist, brauche ich von dir Informationen/Screenshots/Links zur
	      CrypTool-Online-Webseite bzw. zum CTO-Plugin-Framework (z.\,B.
	      Architektur-Dokumentation, bestehende Plugin-Beispiele) -- bitte dann
	      Bescheid geben.
\end{enumerate}
\end{tcolorbox}

\end{document}
